\chapter{Exemplos (eliminar capítulo na versión final)}

\section{Un exemplo de sección}
Esta é {\it letra cursiva}, esta é {\bf letra grosa}, esta é \underline{letra subliñada}, e esta é {\tt letra curier}. Letra {\tiny tiny}, {\scriptsize scriptsize}, {\small small}, {\large large}, {\Large Large}, {\LARGE LARGE} e moitas más. Exemplo de fórmula: $a=\int_o^\infty f(t)dt$.  E agora unha ecuación aparte:

\begin{equation}
    S=\sum_{i=0}^{N-1} a_i^2 .
\label{mi_ecuacion}
\end{equation}

As ecuacións pódense referenciar: ecuación (\ref{mi_ecuacion}).

\subsection{Un exemplo de subsección}
O texto vai aquí.
\subsection{Outro exemplo de subsección}
O texto vai aquí.
\subsubsection{Un exemplo de subsubsección}
O texto vai aquí.
\subsubsection{Un exemplo de subsubsección}
O texto vai aquí.
\subsubsection{Un exemplo de subsubsección}
O texto vai aquí.

\section{Exemplos de figuras e táboas}

A Figura \ref{fig:figura_principal} mostra...

A Táboa \ref{tab:resultados} conten os resultados de...


\begin{figure}
    \centering
    \includegraphics[width=0.9\linewidth]{figuras/figura01.eps}
    \caption{Esta é a figura de tal e cal.}
    \label{fig:figura_principal}
\end{figure}


\begin{table}
    \centering
    \begin{tabular}{l|rc} \toprule
            Izquierda & Derecha & Centrado  \\ \midrule
            ll & r & cccc \\
            llll & rrr & c \\ \bottomrule
        \end{tabular}
        \caption{Esta é a táboa de tal e cal.}
    \label{tab:resultados}
\end{table}

\section{Exemplos de referencias á bibliografía}
Este é un exemplo de referencia a un documento descargado da web \cite{cudaguide}. Agora un libro \cite{goodfellow2016deep}, unha referencia a un artigo de revista \cite{lecun2015deep} e un artigo publicado nun congreso \cite{NIPS2017_3f5ee243}. Tamén se poden pór varias referencias á vez \cite{lecun2015deep,goodfellow2016deep}.


\section{Exemplos de enumeracións}

Con puntos:

\begin{itemize}
    \item Un.
    \item Dous.
    \item Tres.
\end{itemize}

Con números:

\begin{enumerate}
    \item Catro.
    \item Cinco.
    \item Seis.
\end{enumerate}

Exemplo de texto verbatim:

\begin{verbatim}
O texto        verbatim 
     visualizase tal e
            como se escribe
\end{verbatim}

Exemplo de código C:

\lstset{language=C}

\begin{lstlisting}
#include <math.h>
main()
{  int i, j, a[10];
   for(i=0;i<=10;i++) a[i]=i; // comentario 1
   if(a[1]==0) j=1; /* comentario 2 */
   else j=2;
}
\end{lstlisting}

Exemplo de código Java:

\lstset{language=java}

\begin{lstlisting}
class HelloWorldApp {
    public static void main(String[] args) {
        System.out.println("Hello World!"); // Display the string.
    }
}
\end{lstlisting}


